\section{Finite-noise certificates that permit expansion and contraction}
\label{sec:numeric-certificates}

Numerical stability need not be enforced by making every layer nonexpansive
in the ordinary Euclidean norm. Diagonal quadratic metrics provide a
global, finite-noise certificate that allows intermediate expansion and
subsequent contraction. They also give a concrete optimization objective.

Write $H=L-1$ for the number of hidden layers. Let $M_s$ be diagonal
positive semidefinite matrices, and suppose
\begin{equation}\label{eq:diagonal-storage}
 M_H\succeq vv^\top,\qquad
 (W^s)^\top M_sW^s\preceq M_{s-1}\quad(2\le s\le H).
\end{equation}
Define
\[
 Q_M=N\sum_{s=1}^H\lambda_{\max}(M_s).
\]

\begin{proposition}[A diagonal-metric finite-noise certificate]
\label{prop:diagonal-numeric}
Under~\eqref{eq:diagonal-storage}, for every input and every finite
simultaneous postactivation and output disturbance,
\[
 |z_w(x;\eta)-z_w(x)|\le\sqrt{1+Q_M}\,\|\eta\|_*.
\]
In particular this is a global finite-radius statement, not a bound only
on the clean Jacobian.
\end{proposition}
\begin{proof}
Coordinatewise $1$-Lipschitzness of $\tanh$ makes it nonexpansive in every
diagonal seminorm $\|u\|_{M_s}=\sqrt{u^\top M_su}$.
If $\Delta h_s$ denotes the difference of the noisy and clean states,
\eqref{eq:diagonal-storage} gives
\[
 \|\Delta h_s\|_{M_s}
 \le\|\Delta h_{s-1}\|_{M_{s-1}}+\|\eta_s\|_{M_s}.
\]
The first layer has $\Delta h_1=\eta_1$, since there is no input noise.
Therefore
\[
 |\Delta z|\le
   \sum_{s=1}^H\sqrt{\lambda_{\max}(M_s)}\,\|\eta_s\|_2
       +|\eta_{\rm out}|.
\]
Cauchy--Schwarz in the RMS-normalized disturbance coordinates gives the
claim. No assumption on activation signs, margins, or smallness of the
perturbations has been used.
\end{proof}

If all internal operator norms are at most one, taking
$M_s=\|v\|^2I$ recovers $Q_M=H E_v$. But the metric condition is weaker.
For example, consider a scalar hidden stack with alternating internal
weights $2,1/2,2,1/2,\ldots$. Let
$p_s=\prod_{j=s+1}^H|W^j|$ and take $M_s=v^2p_s^2$.
Every matrix inequality in~\eqref{eq:diagonal-storage} is an equality,
every $p_s\le2$, and $Q_M\le4H E_v$. The noise gain is therefore only
$O(\sqrt H)$ despite expansive layers at a positive fraction of depths.
The product of the factors $\max\{1,|W^j|\}$ would instead be exponential
in depth. The certificate accounts for the compensating contractions.

The certificate is also output-specific. For instance, with two hidden
coordinates, readout $v=(1,0)^\top$, and every internal matrix
$W^s=\operatorname{diag}(1,10)$, take
$M_s=\operatorname{diag}(1,0)$. The second coordinate may have a very large
local gain but never affects the output; the certificate gives
$Q_M=H E_v$ nonetheless. Thus irrelevant expansive directions need not
be suppressed.

\paragraph{Computing and training the certificate.}
For fixed weights, minimizing $\sum_s b_s$ subject to
\eqref{eq:diagonal-storage} and $0\preceq M_s\preceq b_sI$, with diagonal
$M_s$, is a semidefinite program: every displayed constraint is affine in
the diagonal metric entries and $b_s$. Its value certifies a finite-noise
gain. One can penalize this certified gain during training, or use the
parameterization
\[
 W^s=D_s^{-1}A_sD_{s-1},\qquad \|A_s\|_{\rm op}\le1,
 \qquad v=D_Hu,\quad\|u\|_2\le1,
\]
with positive diagonal $D_s$ and $M_s=D_s^2$.
Spectral normalization or projection enforces the constraints on $A_s$;
the diagonal scales allow expansion and contraction in the original
coordinates. This is a proposed training constraint or regularizer, not
an assertion that unconstrained gradient descent automatically satisfies it.

For the numerical radius~\eqref{eq:numeric-radius}, a certificate satisfying
\[
 \frac{1+Q_M}{1+H E_v}\le\left(\frac{1000}{16}\right)^2
\]
proves the prior's $1/16$ output-stability condition. More generally, let
$R(L)\ge1$ be any fixed polynomial and replace that radius by
\begin{equation}\label{eq:poly-transient-radius}
 \frac{1}{1000R(L)\sqrt{1+H E_v}}.
\end{equation}
A certificate $Q_M\le R(L)^2H E_v$ then guarantees output change at most
$1/1000$ throughout the ball. The polynomial-depth learning theorem remains
valid: the target cylinder still belongs to the support because the radius
is smaller, and the common compression radius after the readout-tail
truncation is still inverse polynomial in depth.
This permits polynomial intermediate amplification, not just contraction.
For example, internal maps with operator norms at most $1+B/L$ have suffix
gain at most $e^B$; if $B=O(\log L)$, this is polynomial and fits such a
calibration. Applying this observation to a different residual-network
architecture would still require matching that architecture's prior and
target-mass calculation.

The metric inequalities are sufficient certificates for membership in the
numerical prior. If one instead restricts the prior to these particular
certificates, the existing digital target cylinder is not automatically
known to remain inside that stricter support; that would be an additional
expressivity and probability question.

\subsection{A finite-tube version that uses saturation}

The global diagonal-metric test can be conservative for saturated units.
There is a rigorous finite-tube variant. Fix an input, a proposed radius
$r$, and positive definite diagonal $M_s$. Let $a_s=W^sh_{s-1}$ be the
clean preactivations. Put
\[
 B_s=N\sum_{j=1}^s\lambda_{\max}(M_j),\qquad
 R_{s,i}=r\sqrt{B_{s-1}}
       \sqrt{(W^sM_{s-1}^{-1}(W^s)^\top)_{ii}}
       \quad(s\ge2),
\]
and define diagonal matrices $U_s$ by
\[
 (U_s)_{ii}
  =\sup_{|q|\le R_{s,i}}\operatorname{sech}^2(a_{s,i}+q)
  =\operatorname{sech}^2\bigl((|a_{s,i}|-R_{s,i})_+\bigr).
\]
Suppose
\begin{equation}\label{eq:tube-metric}
 M_H\succeq vv^\top,\qquad
 (W^s)^\top U_sM_sU_sW^s\preceq M_{s-1}\quad(2\le s\le H).
\end{equation}
Then every disturbance with $\|\eta\|_*\le r$ changes the output at this
input by at most $r\sqrt{1+B_H}$.

Indeed, induction gives
$\|\Delta h_s\|_{M_s}\le\sum_{j\le s}\|\eta_j\|_{M_j}
\le r\sqrt{B_s}$.
The preceding-layer bound places each preactivation displacement in its
specified interval $[-R_{s,i},R_{s,i}]$. The coordinatewise mean-value
slopes are consequently between zero and the corresponding entries of
$U_s$. Since $M_s$ is diagonal,~\eqref{eq:tube-metric} gives the same
contraction inequality as before. This closes the induction and proves the
finite-noise bound, including scalar output noise by Cauchy--Schwarz.
The construction is self-consistent: the interval widths are computed
before accepting the matrix inequalities, rather than assumed from the
clean Jacobian alone.

To certify the prior, this test must hold at every required input, with
the resulting output bound at most $1/16$. Metrics may be chosen separately
at different inputs. Checking only a finite test set supplies evidence on
those inputs, not a global certificate. Unlike the global test, optimizing
the finite-tube metrics is not in general a single convex semidefinite
program, because the slope bounds depend on the candidate metrics.

\subsection{Large output amplification does not automatically repair this condition}

Let $w$ be fixed with $E_v>0$, and let $w^{(A)}$ be obtained by multiplying
only its readout by $A>0$.
For an input $x$, let
\[
 G_{\rm hid}(w,x)^2
  =N\sum_{s=1}^H
       \left\|\frac{\partial z_w(x;\eta)}{\partial\eta_s}
                     \bigg|_{\eta=0}\right\|_2^2.
\]
This is the squared clean gain from RMS-normalized hidden disturbances;
it does not include the scalar output-noise coordinate. With $c=1/1000$
and $r_A=c/\sqrt{1+H A^2E_v}$,
\begin{equation}\label{eq:readout-scaling-limit}
 \lim_{A\to\infty}
   \sup_{\|\eta\|_*\le r_A}
     |z_{w^{(A)}}(x;\eta)-z_{w^{(A)}}(x)|
   =\frac{c\,G_{\rm hid}(w,x)}{\sqrt{H E_v}}.
\end{equation}
The convergence is uniform over the finite Boolean cube.
To prove it, write $\eta=r_Au$. Smoothness gives
\[
 z_{w^{(A)}}(x;r_Au)-z_{w^{(A)}}(x)
   =A r_A\,\ell_x(u_{\rm hid})+r_Au_{\rm out}
                                      +O(A r_A^2),
\]
uniformly on the unit disturbance ball, where $\ell_x$ is the clean hidden
linear response. Since $A r_A\to c/\sqrt{HE_v}$,
$r_A\to0$, and $A r_A^2\to0$, taking the supremum gives
\eqref{eq:readout-scaling-limit}. Finiteness of the cube makes the remainder
uniform over inputs.

Consequently, if
\[
 \max_x\frac{G_{\rm hid}(w,x)}{\sqrt{H E_v}}>62.5,
\]
then no sufficiently large readout amplification makes the network satisfy
the numerical support rule. For a nonexpansive hidden stack the ratio is
at most one. This distinguishes numerical stability from the earlier
calibrated sign condition, whose required radius could always be met by
sufficiently large amplification of a nonvanishing finite-cube network.
The limiting clean gain is used here only for an asymptotic scaling
obstruction; the positive certificates above control finite disturbances.
